\section{术语表}

\begin{longtable}{@{}p{3.4cm}p{10cm}@{}}
\caption{本文相关术语}\\
\toprule
\textbf{术语} & \textbf{解释} \\
\midrule
\endfirsthead
\toprule
\textbf{术语} & \textbf{解释} \\
\midrule
\endhead
Fast follower & 快速跟随者。不是简单复制，而是在已有技术路线被证明后，快速吸收、工程化、扩展和降本。 \\
\addlinespace
Agentic workflow & 让模型参与多步骤任务、工具调用、规划、执行和反馈循环的工作流。 \\
\addlinespace
RL post-training & 预训练后通过强化学习、人类反馈、环境反馈等方式改进模型行为的阶段。 \\
\addlinespace
Technology ownership mentality & 技术所有权心态，指公司倾向掌握关键底层技术，而不是完全依赖外部供应商。 \\
\addlinespace
Open-weight model & 开放权重模型，可下载部署，但训练数据、完整代码和 recipe 未必开放。 \\
\addlinespace
Closed lab & 权重和核心训练细节不开放，主要通过 API 或产品提供能力的实验室。 \\
\addlinespace
Research taste & 研究品味，指判断哪些方向值得做、如何取舍、如何执行的能力。 \\
\addlinespace
Build-not-buy & 能自建就不购买，尤其在关键技术栈上倾向内部掌握。 \\
\addlinespace
Inference demand & 推理需求，即模型被实际调用、服务用户或业务流程的需求。 \\
\addlinespace
Frontier training & 训练最前沿模型所需的大规模训练过程，通常对算力、数据、系统和稳定性要求极高。 \\
\bottomrule
\end{longtable}


